\documentclass[12pt]{article}
\usepackage[T1]{fontenc}
\usepackage{mathptmx}
\usepackage{amsmath}
\usepackage{amssymb}
\usepackage{array}
\usepackage{booktabs}
\usepackage[a4paper,top=1.6cm,bottom=1.6cm,left=2.3cm,right=2.3cm]{geometry}
\setlength{\parindent}{0pt}
\setlength{\parskip}{4pt}

\newcommand{\ehat}[1]{\mathbf{e}_{#1}}
\newcommand{\rhat}{\hat{\mathbf{r}}}
\newcommand{\that}{\hat{\mathbf{t}}}
\newcommand{\nhat}{\hat{\mathbf{n}}}
\newcommand{\dir}{\operatorname{dir}}

\begin{document}

\begin{center}
\large\textbf{Hall-Sensor Placement Geometry on a Conical Pole}\\[2pt]
\normalsize Derivation of $\;\mathbf{p}=\mathbf{T}+s_{\mathrm{ax}}\ehat{2}+R(s_{\mathrm{ax}})\,\rhat+a\,\ehat{3}$
\end{center}

\vspace{2pt}
\hrule
\vspace{8pt}

\textbf{Scope.} This note derives the closed-form construction used by
\texttt{utils/pole\_sensor\_geometry.m}, which is the single source of truth for the
six Hall-sensor centres $\mathbf{p}_i$ and their outward normals $\nhat_i^{+}$.
All vectors are expressed in the WP frame (origin at the common sphere centre of the
six pole tips) in metres. The construction is per-pole; the index $i$ is dropped.

\section{Direction convention}

Every direction below is written with the elevation--azimuth helper
\begin{equation}
\dir(\varepsilon,\theta)=
\begin{bmatrix}\cos\varepsilon\cos\theta\\ \cos\varepsilon\sin\theta\\ \sin\varepsilon\end{bmatrix},
\qquad
\dir(\varepsilon,\theta)\cdot\dir(\varepsilon',\theta)=\cos(\varepsilon-\varepsilon').
\label{eq:dir}
\end{equation}
The second identity is the only algebra needed to check orthogonality inside a
meridian plane: two directions sharing the azimuth $\theta$ are orthogonal iff their
elevations differ by $90^\circ$.

Per-pole inputs: azimuth $\theta$ (pole angles $0,180,120,300,60,240^\circ$ for
$P_1\ldots P_6$) and the layer flag (lower / upper).

\section{Pole tip and pole axis}

\textbf{Tip.} The six tips lie on a sphere of radius $\ell=R_{\mathrm{norm}}$. The
magic-angle elevation follows from the locked hexapole constraints
$R_{\mathrm{norm},xy}=\ell\sqrt{2/3}$, $R_{\mathrm{norm},z}=\ell/\sqrt{3}$:
\begin{equation}
\psi_0=\arctan\!\frac{R_{\mathrm{norm},z}}{R_{\mathrm{norm},xy}}
      =\arctan\frac{1}{\sqrt2}=35.2644^\circ ,
\qquad
\ehat{1}=\dir\!\big(\sigma_\ell\psi_0,\ \theta\big),\qquad
\mathbf{T}=\ell\,\ehat{1},
\label{eq:tip}
\end{equation}
with $\sigma_\ell=-1$ for a lower pole and $+1$ for an upper pole.

\textbf{Axis.} The cone axis $\ehat{2}$ points from the tip \emph{towards the base}
(direction of increasing axial station $s$):
\begin{equation}
\ehat{2}=\dir(\varepsilon,\theta),\qquad
\varepsilon=\begin{cases}
0 & \text{lower pole (horizontal axis)}\\
\varepsilon_{\mathrm{up}}=36.5895^\circ & \text{upper pole}.
\end{cases}
\label{eq:axis}
\end{equation}

\textbf{Remark (do not confuse $\ehat{1}$ with $\ehat{2}$).} The tip direction and the
pole axis are \emph{different} vectors. For a lower pole $\ehat{1}$ has elevation
$-35.2644^\circ$ while $\ehat{2}$ is horizontal; for an upper pole the elevations are
$+35.2644^\circ$ and $+36.5895^\circ$. Only $\ehat{1}$ carries the magic angle --- it is
fixed by the orthogonality/common-sphere constraints. The axis inclination
$\varepsilon_{\mathrm{up}}$ is the physical tilt of the machined pole and is an
independent number.

\section{The local triad and the radial vector \boldmath$\rhat$}

This is the part that is easy to forget. Fix a pole. Two more unit vectors are built
so that $\{\ehat{2},\rhat_0,\that\}$ is orthonormal:
\begin{align}
\rhat_0 &= \dir\!\big(\varepsilon+\sigma\,90^\circ,\ \theta\big),
  &&\sigma=\begin{cases}-1 & \text{lower pole}\\ +1 & \text{upper pole}\end{cases}
  \label{eq:rad0}\\[2pt]
\that &= \begin{bmatrix}-\sin\theta\\ \cos\theta\\ 0\end{bmatrix}.
  \label{eq:tanv}
\end{align}

\textbf{Why these are orthonormal.} $\ehat{2}$ and $\rhat_0$ share the azimuth
$\theta$ and their elevations differ by exactly $90^\circ$, so
$\ehat{2}\cdot\rhat_0=\cos(\mp90^\circ)=0$ by \eqref{eq:dir}. The vector $\that$ is
horizontal and perpendicular to the azimuthal direction
$[\cos\theta,\sin\theta,0]^{\mathsf T}$, hence perpendicular to \emph{any}
$\dir(\cdot,\theta)$; in particular $\that\perp\ehat{2}$ and $\that\perp\rhat_0$.

\textbf{Geometric meaning.}
\begin{itemize}\itemsep2pt
\item $\ehat{2},\rhat_0$ span the \emph{meridian plane} --- the vertical plane
      containing the pole axis (azimuth $\theta$). $\rhat_0$ is the radial direction
      inside that plane, pointing away from the axis.
\item The sign $\sigma$ selects \emph{which side} of the axis the sensor is mounted on.
      A lower pole is the half-cut part: the mounting face is its \textbf{bottom} cone
      surface, so $\sigma=-1$ tilts $\rhat_0$ downwards. An upper pole is a full cone
      and is mounted on the upper side, so $\sigma=+1$.
\item $\that$ is tangential --- perpendicular to the meridian plane. It is only needed
      to rotate the mounting point around the pole axis.
\end{itemize}

\textbf{Azimuthal placement around the axis.} Rotating $\rhat_0$ about $\ehat{2}$ by an
angle $\psi$ stays in the plane spanned by $\rhat_0$ and $\that$ (both are
perpendicular to $\ehat{2}$), so Rodrigues' formula collapses to
\begin{equation}
\boxed{\ \rhat(\psi)=\cos\psi\;\rhat_0+\sin\psi\;\that\ }
\label{eq:rhat}
\end{equation}
with $\rhat(\psi)\perp\ehat{2}$ and $\|\rhat(\psi)\|=1$ for all $\psi$. The canonical
sensor azimuth is $\psi=0$, i.e. $\rhat=\rhat_0$, which places the sensor in the
meridian plane --- directly below (lower) or above (upper) the pole axis.

\textbf{Caveat on the sense of rotation.} A direct evaluation of the cross product
gives $\ehat{2}\times\rhat_0=-\sigma\,\that$. Therefore \eqref{eq:rhat} is a
right-handed rotation about $\ehat{2}$ by $-\sigma\psi$: the sense of increasing $\psi$
is opposite between the two layers. This is immaterial at the default $\psi=0$ and for
symmetric sweeps, but it must be kept in mind if a signed azimuthal offset is ever
imposed.

\section{The real cone: virtual apex, radius, and outward normal}

\subsection{Virtual apex}

The pole is not a mathematically sharp cone: the tip is rounded with a fillet of
radius $r_f=40\ \mu$m, and the conical surface is \emph{tangent} to that sphere. Let
$\beta$ be the true half-cone angle. For a sphere of radius $r_f$ inscribed tangentially
in a cone of half-angle $\beta$, the distance from the apex to the sphere centre is
$r_f/\sin\beta$. Placing the sphere centre at axial station $s=r_f$ (so that the
foremost point of the sphere --- the physical tip $\mathbf{T}$ --- sits at $s=0$), the
apex lies \emph{outside} the material at
\begin{equation}
s=-e_v,\qquad
\boxed{\ e_v=\frac{r_f}{\sin\beta}-r_f\ }\;\approx\;0.16\ \mathrm{mm}.
\label{eq:apex}
\end{equation}

\subsection{Cone radius}

Measuring the axial station $s$ from the physical tip, the distance from the virtual
apex is $e_v+s$, so
\begin{equation}
R(s)=(e_v+s)\tan\beta,
\qquad
R(0)=e_v\tan\beta\approx 0.033\ \mathrm{mm}\neq 0 .
\label{eq:radius}
\end{equation}
The non-zero intercept is the whole point: \textbf{at the axial position of the tip the
cone surface already has a finite radius}. Anchoring the mounting point at the ideal
tip with $R=0$ is what used to push the sensor into the steel (Sec.~\ref{sec:old}).

\subsection{Outward normal}

Parametrise the cone surface by $(s,\psi)$:
\begin{equation}
\mathbf{P}(s,\psi)=\mathbf{T}+s\,\ehat{2}+R(s)\,\rhat(\psi).
\end{equation}
The generator (slant) direction is
\begin{equation}
\hat{\mathbf{g}}
=\frac{\partial\mathbf{P}/\partial s}{\|\partial\mathbf{P}/\partial s\|}
=\frac{\ehat{2}+\tan\beta\,\rhat}{\sqrt{1+\tan^2\beta}}
=\cos\beta\,\ehat{2}+\sin\beta\,\rhat .
\label{eq:gen}
\end{equation}
The surface normal lies in the meridian-like plane $\mathrm{span}\{\ehat{2},\rhat\}$,
is perpendicular to $\hat{\mathbf{g}}$, and must point away from the axis
($\nhat\cdot\rhat>0$). The unique unit vector satisfying all three is
\begin{equation}
\boxed{\ \ehat{3}\;=\;\nhat^{+}\;=\;\cos\beta\,\rhat-\sin\beta\,\ehat{2}\ }
\label{eq:normal}
\end{equation}
since $\hat{\mathbf{g}}\cdot\ehat{3}=\cos\beta\sin\beta-\sin\beta\cos\beta=0$ and
$\ehat{3}\cdot\rhat=\cos\beta>0$. At $\psi=0$ this is exactly
$\dir\!\big(\varepsilon+\sigma(\beta+90^\circ),\theta\big)$, i.e. the axis direction
rotated by $\sigma(90^\circ+\beta)$ inside the meridian plane.

\section{Axial station from the slant offset}

$\mathrm{SOFF}$ is defined as the \textbf{slant distance measured along the generator},
starting from the ``red point'' --- the point of the generator at the same axial station
as the tip, i.e. $(s,R)=(0,\,e_v\tan\beta)$. Advancing a slant distance
$\mathrm{SOFF}$ along $\hat{\mathbf{g}}$ advances the axial station by its $\ehat{2}$
component, hence
\begin{equation}
\boxed{\ s_{\mathrm{ax}}=\mathrm{SOFF}\cdot\cos\beta\ }
\label{eq:sax}
\end{equation}
with the design value $\mathrm{SOFF}=4.572$~mm ($=0.18$~inch).

\textbf{Convention.} Quoting a mounting offset always means the \emph{full} slant
distance measured in the meridian plane, \emph{including} its leading segment: saying
``$\mathrm{SOFF}=4$~mm'' already contains the portion $t_{\tan}$ that is pressed
against the fillet rather than the cone, and the whole $4$~mm is then projected onto
the axis by $\cos\beta$. The CAD-measured fillet segments are $t_{\tan}=0.0328$~mm
(lower) and $0.0311$~mm (upper); they are recorded only for drawing the pole outline
and do \emph{not} enter \eqref{eq:sax}.

\textbf{History.} An earlier revision used
$s_{\mathrm{ax}}=(\mathrm{SOFF}-t_{\tan})\cos\beta+t_{\tan}$, which measured
$\mathrm{SOFF}$ from the physical tip and treated the leading segment as purely
axial. The two differ by $t_{\tan}(1-\cos\beta)\approx0.6\ \mu$m, three orders of
magnitude below the air gap, so no downstream result changed.

\section{Assembly}

Stacking \eqref{eq:tip}, \eqref{eq:sax}, \eqref{eq:radius}, \eqref{eq:rhat} and
\eqref{eq:normal}, and lifting the sensor off the surface by the air gap
$a=0.41$~mm along the outward normal:
\begin{equation}
\boxed{\;
\mathbf{p}
=\underbrace{\mathbf{T}}_{\text{tip}}
+\underbrace{s_{\mathrm{ax}}\,\ehat{2}}_{\text{slide along axis}}
+\underbrace{R(s_{\mathrm{ax}})\,\rhat(\psi)}_{\text{out to the cone surface}}
+\underbrace{a\,\ehat{3}}_{\text{lift off the surface}}\;}
\label{eq:assembly}
\end{equation}

Read left to right, the four terms are: start at the pole tip; slide along the axis to
the correct axial station; step radially outwards until you reach the actual cone
surface at that station; then lift perpendicular to that surface by the air gap.
Because the last step uses the true surface normal \eqref{eq:normal}, the perpendicular
distance from the sensor centre to the cone surface is \emph{exactly} $a$ by
construction.

\subsection*{Degenerate case: flat mounting face}

For the lower pole the option \texttt{face\_lower='flat'} mounts the sensor on the
half-cut plane instead of the bottom cone. That plane \emph{contains} the pole axis,
so there is no radial stand-off and no slant projection:
\begin{equation}
s_{\mathrm{ax}}=\mathrm{SOFF},\qquad R=0,\qquad \nhat^{+}=\hat{\mathbf{z}},
\qquad \mathbf{p}=\mathbf{T}+\mathrm{SOFF}\,\ehat{2}+a\,\hat{\mathbf{z}} .
\end{equation}
The azimuth $\psi$ has no meaning for a plane and is ignored.

\section{Numerical values (long2016 half-cut, CAD-measured cone)}
\label{sec:num}

$\ell=0.5$~mm, $r_f=0.04$~mm, $\mathrm{SOFF}=4.572$~mm, $a=0.41$~mm.
The two layers have \emph{different} measured cone slopes
($\tan\beta=0.20330$ lower, $0.19463$ upper; a $4.4\%$ difference).

\begin{center}
\begin{tabular}{lccccc}
\toprule
Layer & $\beta$ [deg] & $e_v$ [mm] & $s_{\mathrm{ax}}$ [mm] & $R(s_{\mathrm{ax}})$ [mm] & poles \\
\midrule
Lower & 11.4916 & 0.16078 & 4.48035 & 0.94354 & $P_1,P_3,P_6$ \\
Upper & 11.0138 & 0.16937 & 4.48779 & 0.90642 & $P_2,P_4,P_5$ \\
\bottomrule
\end{tabular}
\end{center}

Spot check for $P_1$ ($\theta=0$, lower, $\psi=0$): $\ehat{2}=[1,0,0]^{\mathsf T}$,
$\rhat_0=\dir(-90^\circ,0)=[0,0,-1]^{\mathsf T}$, so \eqref{eq:normal} gives
\[
\nhat^{+}=\cos\beta\,[0,0,-1]^{\mathsf T}-\sin\beta\,[1,0,0]^{\mathsf T}
=[-0.1992,\ 0,\ -0.9800]^{\mathsf T},
\]
which is what the code returns. The normal tilts $\beta$ away from straight-down,
as it must for a surface whose generator is inclined by $\beta$ to the axis.

\section{Why the naive construction fails}
\label{sec:old}

A tempting shortcut is to anchor at the ideal tip and walk the slant distance directly:
$\mathbf{p}=\mathbf{T}+\mathrm{SOFF}\cdot\dir(\varepsilon+\sigma\beta,\theta)+a\,\ehat{3}$.
This ignores the finite radius $R(0)=e_v\tan\beta\approx0.033$~mm at the tip station,
so the mounting point is placed \emph{inside} the steel. The resulting perpendicular
distance to the \emph{real} cone surface is
\[
0.3635\ \mathrm{mm}\ \text{(lower)},\qquad 0.4012\ \mathrm{mm}\ \text{(upper)},
\]
instead of the specified $0.41$~mm --- an $11\%$ error on the lower poles, and
inconsistent between layers. The construction \eqref{eq:assembly} removes both defects:
the stand-off is exactly $a$ on every pole.

\vspace{6pt}
\hrule
\vspace{4pt}
{\footnotesize
\textbf{Implementation:} \texttt{matlab/Maxwell/utils/pole\_sensor\_geometry.m}
(and the identical copy under \texttt{matlab/APDL/Calibration\_using\_FEM\_modeling/utils/}).
Consumers: \texttt{build\_V\_matrix.m} and every script that draws sensors or pole
outlines --- none of them may re-derive this geometry locally.\par
\textbf{Related:} \texttt{.claude/rules/ansys-cad-alignment.md} (CAD is the source of
truth); memory \texttt{long2016-pole-cone-geometry} (STEP measurements and how they
were taken).}

\end{document}
